\section{Prácticas recomendadas para el diseño de herramientas}

El Apéndice II del artículo es un capítulo extremadamente valioso desde una perspectiva práctica, que discute cómo diseñar interfaces de herramientas para los Agentes; este es un factor determinante para el rendimiento de los Agentes.

\subsection{Interfaz Agente-Computadora (ACI)}

Anthropic propuso un concepto importante: \textbf{ACI} (Agent-Computer Interface), que corresponde a HCI (Human-Computer Interface). La idea central es: al diseñar interfaces de herramientas para Agentes, se debe invertir la misma energía que en el diseño de interfaces humano-computadora.

\begin{importantbox}{Principios de diseño de ACI}
\begin{enumerate}
    \item \textbf{Pensamiento empático}: Desde la perspectiva del modelo, ¿se puede saber intuitivamente cómo usar esta herramienta solo con la descripción y los parámetros? Si usted mismo necesita reflexionar cuidadosamente, el modelo también lo hará.
    \item \textbf{Documentación completa}: Una buena definición de herramientas debe incluir ejemplos de uso, casos límite, requisitos de formato de entrada y diferencias con otras herramientas.
    \item \textbf{Pruebas iterativas}: Pruebe en Workbench cómo el modelo usa las herramientas con grandes volúmenes de entradas de ejemplo e itere para optimizar según los errores observados.
    \item \textbf{Diseño a prueba de errores (Poka-yoke)}: Modifique el diseño de parámetros para hacer más difícil que el modelo cometa errores.
\end{enumerate}
\end{importantbox}

\subsection{Selección del formato de herramientas}

Una misma operación suele tener múltiples formas de descripción. Por ejemplo, la edición de archivos puede usar el formato diff o reescribir todo el archivo; la salida estructurada puede ser Markdown o JSON. Aunque en ingeniería de software estas diferencias son superficiales (se pueden convertir sin pérdida), para un LLM las diferentes formatos presentan \textbf{diferencias enormes en dificultad}.

\begin{table}[H]
\centering
\begin{tabular}{>{\bfseries}p{3cm} p{4.5cm} p{4.5cm}}
\toprule
Dimensión de diseño & Práctica recomendada & Práctica a evitar \\
\midrule
Espacio de razonamiento & Proveer al modelo suficientes tokens para «pensar» antes de escribir la salida clave & Exigir que el modelo determine metainformación (como el número de líneas del diff) antes de escribir el contenido \\
Naturalidad del formato & Usar formatos comunes en textos de internet para los modelos & Usar formatos personalizados poco frecuentes o altamente estructurados \\
Costo del formato & Minimizar la carga del formato & Exigir conteos exactos de miles de líneas de código o realizar escape JSON del código \\
Especificación de ruta & Usar rutas absolutas & Usar rutas relativas (prone a errores por cambios de directorio) \\
\bottomrule
\end{tabular}
\caption{Prácticas recomendadas y prácticas a evitar en el diseño de formatos de herramientas}
\end{table}

\begin{knowledgebox}{Experiencia práctica proveniente de SWE-bench}
Cuando Anthropic construyó el Agente para SWE-bench, \textbf{el tiempo invertido en optimización de herramientas superó con mucho la optimización general del prompt}. Un caso típico: el modelo comete errores frecuentes al usar herramientas con rutas relativas (porque el Agente cambia el directorio de trabajo); al cambiar las herramientas para exigir rutas absolutas obligatoriamente, el uso del modelo resultó perfecto. Esto demuestra que un buen diseño de ACI puede eliminar desde la raíz ciertas clases de errores.
\end{knowledgebox}

\subsection{Resumen de la sección}
El diseño de herramientas es un factor clave para el éxito de los Agentes, incluso más importante que la optimización del prompt. El diseño de ACI debe seguir los cuatro principios: pensamiento empático, documentación completa, pruebas iterativas y diseño a prueba de errores. Se deben seleccionar formatos de herramienta naturales que sean fáciles de generar por el modelo, evitando cargas de formato innecesarias.
